\section{Bayes' theorem}

The law of total probability lets us compute the probability of an observation
by splitting the world into possible cases. Bayes' theorem answers the reverse
question. After an observation is known, how should we update the probabilities
of the possible cases?

This is not a separate trick. It is conditional probability applied in the
reverse direction.

\subsection*{The basic reversal problem}

Suppose \(B\) is a possible background case and \(A\) is an observed event. We
may know or be able to model
\[
P(A\mid B),
\]
the probability of observing \(A\) if case \(B\) is true. But after observing
\(A\), we may want
\[
P(B\mid A),
\]
the probability that case \(B\) is true given that \(A\) was observed.

These two probabilities are not the same. They use different reference spaces.
The first reasons inside \(B\). The second reasons inside \(A\).

From the multiplication rule,
\[
P(A\cap B)=P(A\mid B)P(B).
\]
Also,
\[
P(A\cap B)=P(B\mid A)P(A).
\]
Since both expressions describe the same intersection, we get
\[
P(B\mid A)P(A)=P(A\mid B)P(B).
\]
If \(P(A)>0\), then
\[
P(B\mid A)=\frac{P(A\mid B)P(B)}{P(A)}.
\]

\begin{theorembox}
\textbf{Bayes' theorem, two-event form.} If \(P(A)>0\) and \(P(B)>0\), then
\[
P(B\mid A)=\frac{P(A\mid B)P(B)}{P(A)}.
\]
\end{theorembox}

This formula says that the updated probability of \(B\) after observing \(A\) is
proportional to two quantities:
\begin{itemize}
    \item how plausible \(B\) was before observing \(A\), namely \(P(B)\);
    \item how well \(B\) explains the observation \(A\), namely \(P(A\mid B)\).
\end{itemize}

\subsection*{Why the denominator is often expanded}

In many applications, \(P(A)\) is not known directly. Instead, the observation
\(A\) can occur under several possible cases. If \(B_1,
\ldots,B_n\) form a
partition of \(\Omega\), then the law of total probability gives
\[
P(A)=\sum_{i=1}^n P(A\mid B_i)P(B_i).
\]
Substituting this into Bayes' theorem gives
\[
P(B_j\mid A)=
\frac{P(A\mid B_j)P(B_j)}
{\sum_{i=1}^n P(A\mid B_i)P(B_i)}.
\]

\begin{theorembox}
\textbf{Bayes' theorem over a partition.} If \(B_1,
\ldots,B_n\) is a partition
of \(\Omega\), \(P(B_i)>0\), and \(P(A)>0\), then
\[
P(B_j\mid A)=
\frac{P(A\mid B_j)P(B_j)}
{\sum_{i=1}^n P(A\mid B_i)P(B_i)}.
\]
\end{theorembox}

The denominator is not an arbitrary normalizing term. It is the total probability
of the observation \(A\), obtained by adding all the ways in which \(A\) can
occur through the cases \(B_i\).

\subsection*{Prior, likelihood, posterior}

Bayes' theorem is often described using three words.

The number
\[
P(B_j)
\]
is the \textbf{prior probability} of case \(B_j\). It is the probability before
observing \(A\).

The number
\[
P(A\mid B_j)
\]
is the \textbf{likelihood} of the observation under case \(B_j\). It measures how
compatible the observation is with that case.

The number
\[
P(B_j\mid A)
\]
is the \textbf{posterior probability}. It is the updated probability of the case
after observing \(A\).

The logic is:
\[
\text{prior}\times\text{likelihood}\quad\longrightarrow\quad\text{posterior after normalization}.
\]
The normalization is necessary because the posterior probabilities over all
cases must add to one.

\subsection*{Example: choosing an urn after observing a colour}

Suppose one of two urns is chosen. Urn 1 is chosen with probability \(2/3\), and
urn 2 is chosen with probability \(1/3\). Urn 1 contains \(3\) red objects and
\(1\) blue object. Urn 2 contains \(1\) red object and \(4\) blue objects. One
object is drawn from the chosen urn, and it is red. Which urn is now more likely
to have been chosen?

Let
\[
B_1=\{\text{urn 1 was chosen}\},\qquad B_2=\{\text{urn 2 was chosen}\},
\]
and
\[
A=\{\text{the drawn object is red}\}.
\]
We have
\[
P(B_1)=\frac23,\qquad P(B_2)=\frac13,
\]
\[
P(A\mid B_1)=\frac34,\qquad P(A\mid B_2)=\frac15.
\]
First compute the total probability of drawing red:
\[
P(A)=P(A\mid B_1)P(B_1)+P(A\mid B_2)P(B_2).
\]
Thus
\[
P(A)=\frac34\cdot\frac23+\frac15\cdot\frac13
=\frac12+\frac1{15}=\frac{17}{30}.
\]
Now Bayes' theorem gives
\[
P(B_1\mid A)=\frac{P(A\mid B_1)P(B_1)}{P(A)}.
\]
Therefore
\[
P(B_1\mid A)=\frac{\frac34\cdot\frac23}{\frac{17}{30}}
=\frac{\frac12}{\frac{17}{30}}
=\frac{15}{17}.
\]
Similarly,
\[
P(B_2\mid A)=\frac{\frac15\cdot\frac13}{\frac{17}{30}}
=\frac{\frac1{15}}{\frac{17}{30}}
=\frac{2}{17}.
\]
The posterior probabilities add to one:
\[
\frac{15}{17}+\frac{2}{17}=1.
\]

The red observation strongly supports urn 1 because urn 1 was already more
likely before the draw and also has a higher red proportion.

\subsection*{Example: diagnostic testing}

Suppose a condition occurs in \(1\%\) of a population:
\[
P(C)=0.01,\qquad P(C^c)=0.99.
\]
A test has sensitivity \(95\%\), meaning
\[
P(T\mid C)=0.95,
\]
and false positive rate \(5\%\), meaning
\[
P(T\mid C^c)=0.05.
\]
Here \(T\) is the event that the test result is positive.

We want
\[
P(C\mid T),
\]
the probability that a person has the condition after receiving a positive test
result.

By Bayes' theorem over the partition \(C,C^c\),
\[
P(C\mid T)=
\frac{P(T\mid C)P(C)}{P(T\mid C)P(C)+P(T\mid C^c)P(C^c)}.
\]
Substitute the values:
\[
P(C\mid T)=
\frac{0.95\cdot0.01}{0.95\cdot0.01+0.05\cdot0.99}.
\]
The numerator is
\[
0.0095.
\]
The denominator is
\[
0.0095+0.0495=0.059.
\]
Therefore
\[
P(C\mid T)=\frac{0.0095}{0.059}\approx0.161.
\]
So the probability is about \(16.1\%\), not \(95\%\).

This example is important because it shows the role of the base rate. A positive
test result is evidence for the condition, but when the condition is rare, false
positives may still be numerous among all positive results.

\subsection*{Natural frequency interpretation}

The diagnostic example becomes clearer if we imagine \(10000\) people.

About \(1\%\) have the condition, so approximately
\[
100
\]
people have it, and
\[
9900
\]
do not.

Among the \(100\) people with the condition, the test is positive for about
\[
0.95\cdot100=95
\]
people.

Among the \(9900\) people without the condition, the test is positive for about
\[
0.05\cdot9900=495
\]
people.

So the total number of positive tests is about
\[
95+495=590.
\]
Among these, only \(95\) are true positives. Hence
\[
P(C\mid T)\approx\frac{95}{590}\approx0.161.
\]
This is the same calculation as Bayes' theorem, but written as counts. It often
makes the denominator easier to understand: the positive tests come from two
sources, true positives and false positives.

\subsection*{Bayes' theorem as redistribution of probability over cases}

Suppose \(B_1,
\ldots,B_n\) are possible cases. Before observing \(A\), the
probability assigned to case \(B_i\) is \(P(B_i)\). After observing \(A\), each
case receives support proportional to
\[
P(A\mid B_i)P(B_i).
\]
The product has a clear meaning:
\[
P(A\mid B_i)P(B_i)=P(A\cap B_i).
\]
It is the probability that case \(B_i\) occurs and produces the observation
\(A\).

After observing \(A\), all posterior probabilities must add to one. Therefore we
divide each contribution by the total probability of \(A\):
\[
P(A)=\sum_{i=1}^n P(A\mid B_i)P(B_i).
\]
This turns the contributions into posterior probabilities.

\subsection*{Common reversal error}

The most common mistake is to confuse
\[
P(A\mid B)
\]
with
\[
P(B\mid A).
\]
For example, test sensitivity \(P(T\mid C)\) is not the same as the probability
of having the condition after a positive test \(P(C\mid T)\). The first uses the
condition group as the reference space. The second uses the positive-test group
as the reference space.

Bayes' theorem exists precisely because these two directions are different but
related through the intersection \(A\cap B\).

\begin{remarkbox}
Whenever a conditional probability appears, read the phrase after the vertical
bar first. It tells you the reference space. Then read the event before the bar.
It tells you what is being measured inside that reference space.
\end{remarkbox}

\subsection*{What Bayes' theorem teaches}

Bayes' theorem teaches how probability is updated by evidence. It combines three
ingredients:
\begin{itemize}
    \item prior probabilities of the possible cases;
    \item likelihoods of the observation under those cases;
    \item normalization by the total probability of the observation.
\end{itemize}

\begin{theorembox}
Bayes' theorem is conditional probability used backward. It converts
\[
P(\text{observation}\mid\text{case})
\]
into
\[
P(\text{case}\mid\text{observation}),
\]
using the prior probabilities of the cases and the total probability of the
observation.
\end{theorembox}
